% FILE: 07_open_questions_and_next_work.tex
% PROJECT: ma — M&A Claim Taxonomy and Board Projection Surface

\section{Open Questions and Next Work}
\label{sec:ma-open-questions}

\subsection{Known Gaps}
\label{sec:ma-gaps-known}

Four structural gaps are identified as of the thesis writing date (2026-06-01). These are not
ambiguities in the doctrine---the doctrine is APPROVED. They are gaps between the specified
execution capability and the demonstrated execution capability.

\paragraph{Gap 1: No Real Enterprise Event Logs.}
The five receipt JSON files contain realistic and structurally correct values. The EBITDA
receipt reports fitness 0.982, precision 0.945, and a \$1,250,000 impact. These values are
consistent with the mathematical formulas and admissibility thresholds. However, they were not
produced by executing wasm4pm queries against actual enterprise event logs extracted from a
real ERP or CRM system. No real transaction data room has been connected to this framework.
The receipts are specification-illustrative, not execution-derived. (Source: evidence
extraction from \texttt{receipts/} directory and module structure audit.)

\paragraph{Gap 2: wasm4pm Query Artifacts Unconfirmed on Disk.}
The receipt files reference wasm4pm query URIs at absolute paths (e.g.,
\texttt{file:///Users/sac/process-intelligence/sources/wasm4pm/queries/ebitda\_rework\_reduction.wasm}).
These paths are cited in the receipt \texttt{query\_definition} fields but have not been
confirmed to exist as compiled \texttt{.wasm} binaries. The connection between the claim
taxonomy specifications and the wasm4pm execution engine is asserted by reference but not
demonstrated from within this module.

\paragraph{Gap 3: Reverse Porter Five Forces---No External Adoption Evidence.}
The Reverse Porter Five Forces document is the strongest conceptual differentiator in the
module. Force 4 (substitute prevention via receipt dependency) argues that buyers who build
analytical workflows on wasm4pm receipts cannot switch to a process-evidence competitor
without re-running all historical conformance queries. This is a network-effect claim whose
validity depends on external adoption. As of 2026-06-01, no external buyer or auditor has
adopted wasm4pm receipts. The network effect is theoretical. (AUTHOR THESIS: this is an
honest assessment of the current state, not a defect in the framework's logic.)

\paragraph{Gap 4: HTR-to-Chatman-Equation Lineage Unspecified Within Module.}
The Hostile Takeover Receipt taxonomy (v30.1.1) references ``Otar Generative Pipeline
receipts'' as the basis for unforgeable process replay. The connection between the HTR
taxonomy, the Chatman Equation $A = \mu(O^*)$, and the 2016 language-model experiment
lineage is not made explicit within the \texttt{ma/} module itself. The capital-flow and
AI XYNZ connections to this M\&A surface remain unspecified in module documentation.
(Source: structural audit of \texttt{hostile-takeover-receipts.md}.)

\subsection{Deferred Gates}
\label{sec:ma-gaps-deferred}

The following proof gates are explicitly deferred to the execution phase:

\begin{itemize}
  \item \textbf{Claims data population:} A JSON array of board claims derived from an actual
    transaction data room must be provided as input to the Board Projection Renderer. No
    target company has been identified or engaged.
  \item \textbf{Receipt generation from live wasm4pm execution:} Five receipts must be
    produced by executing wasm4pm queries against real XES/OCEL 2.0 logs and verifying that
    the output matches the specification receipt schema.
  \item \textbf{End-to-end renderer execution:} The seven-step pipeline (Claims JSON $\to$
    Tera2 $\to$ pptx-rs $\to$ calamine $\to$ validation $\to$ ledger seal) must be executed
    to produce a sealed PowerPoint deck and Excel workbook. Rendering libraries are specified
    but not invoked.
  \item \textbf{Auditor path live trial:} The five-step auditor verification protocol must
    be executed by an independent party using only the VDR contents. No such trial has
    occurred.
  \item \textbf{Negative testing:} Non-conforming traces must be injected into wasm4pm
    queries to verify correct rejection. This is specified in the admissibility framework but
    not executed.
\end{itemize}

\subsection{Research Risks}
\label{sec:ma-gaps-risks}

\paragraph{Risk 1: Admissibility Threshold Calibration.}
The thresholds $f \ge 0.95$, $p \ge 0.90$ are grounded in academic process mining literature
(Adriansyah 2014). Their applicability to M\&A due-diligence contexts---where event logs may
be incomplete, noisy, or extracted under adversarial conditions---has not been empirically
validated. Real enterprise event logs frequently exhibit lower fitness due to exception
handling, manual interventions, and system migration artifacts. A target whose best-faith
process achieves fitness 0.88 would fail the T3 admissibility gate as currently specified.
The threshold calibration may require empirical adjustment once real logs are available.

\paragraph{Risk 2: OCEL 2.0 Adoption Lag.}
The module mandates OCEL 2.0 as the preferred event log format for multi-object processes.
As of 2026, OCEL 2.0 tooling adoption in enterprise ERP systems (SAP, Oracle, Salesforce) is
nascent. A real diligence engagement may require XES extraction with manual object-centric
enrichment---a labor-intensive process that partly offsets the 7-to-3-day compression claim.

\paragraph{Risk 3: wasm4pm Query Portability.}
The receipt schema references wasm4pm query URIs as absolute local paths. In a real VDR
context, queries must be portable across buyer and auditor environments. The portability of
compiled \texttt{.wasm} binaries across runtime environments, and the specification of
the wasm4pm runtime interface that auditors must install, are not fully specified in the
current module.

\subsection{Next Receipted Work}
\label{sec:ma-gaps-next}

The following work items are required to advance the module from PARTIAL to ALIVE status.
Each must produce a receipt or checkpoint before the gate is considered closed:

\begin{enumerate}
  \item \textbf{Confirm wasm4pm query artifacts:} Verify that the five query \texttt{.wasm}
    files referenced in receipts exist on disk and execute correctly against synthetic event
    logs. Produce a receipt for each successful execution.
  \item \textbf{Synthetic event log construction:} Manufacture XES and OCEL 2.0 event logs
    with known properties (fitness, entropy, rework rate) to serve as the canonical test
    inputs for the framework. Receipt each log with its BLAKE3 hash.
  \item \textbf{End-to-end renderer execution on synthetic data:} Execute the seven-step
    Board Projection Renderer pipeline using synthetic claims data and synthetic receipts.
    Produce a sealed PowerPoint and Excel artifact. Receipt the ledger root hash.
  \item \textbf{Auditor protocol dry run:} Execute the five-step auditor verification
    protocol against the synthetic VDR. Document which steps pass, which require tooling
    not yet specified, and what the independent replay deviation is.
  \item \textbf{Negative testing:} Inject five classes of non-conforming traces (skipped
    steps, wrong actor, out-of-order events, timestamp violations, SoD violations) and
    verify that the wasm4pm queries and admissibility protocol correctly reject each.
  \item \textbf{HTR-to-Chatman lineage documentation:} Add an explicit section to
    \texttt{hostile-takeover-receipts.md} tracing the HTR taxonomy to the Chatman Equation
    $A = \mu(O^*)$ and the 2016 language-model lineage, establishing the capital-flow
    surface as a named terminus in the thesis architecture.
\end{enumerate}

None of these items alter the doctrine. All of them are execution-layer tasks whose
receipts will constitute the evidence base for elevating the module's status from PARTIAL
to ALIVE.
